% =====================================================================
%  CAPÍTULO 4 — LEVANTAMENTO E ANÁLISE DE REQUISITOS
% =====================================================================
\chapter{Levantamento e Análise de Requisitos}
\label{cap:requisitos}

Este capítulo formaliza os requisitos do sistema, distinguindo os
\emph{funcionais} (o que o sistema deve fazer) dos \emph{não funcionais}
(restrições de qualidade), e descreve os principais casos de uso através
de um diagrama e de \emph{user stories}.

\section{Requisitos Funcionais}
\label{sec:rf}

Os requisitos funcionais derivam diretamente dos objetivos
(\Cref{sec:objetivos}) e do âmbito definido com a orientação. A
\Cref{tab:rf} lista-os, atribuindo a cada um uma prioridade
(\emph{Must}, \emph{Should} ou \emph{Could}, segundo a notação MoSCoW).

\begin{longtable}{@{}p{1.4cm}p{8.8cm}p{2.4cm}@{}}
\caption{Requisitos funcionais do sistema.}\label{tab:rf}\\
\toprule
\textbf{ID} & \textbf{Descrição} & \textbf{Prioridade} \\
\midrule
\endfirsthead
\multicolumn{3}{c}{\tablename\ \thetable\ -- continuação}\\
\toprule
\textbf{ID} & \textbf{Descrição} & \textbf{Prioridade} \\
\midrule
\endhead
\rf{1}  & O sistema deve permitir registo, início e fim de sessão, e
recuperação/alteração de palavra-passe. & Must \\
\rf{2}  & Cada utilizador deve aceder exclusivamente aos seus próprios
dados (isolamento multi-inquilino). & Must \\
\rf{3}  & O sistema deve permitir o carregamento de fotografias,
individualmente e em lote, com \emph{drag-and-drop}. & Must \\
\rf{4}  & O sistema deve importar árvores genealógicas em formato GEDCOM e
construir a estrutura familiar. & Must \\
\rf{5}  & O sistema deve extrair automaticamente metadados \textsc{exif}
das fotografias (data, GPS, câmara). & Must \\
\rf{6}  & O sistema deve descrever o conteúdo visual de cada fotografia
através de um modelo de IA de visão. & Must \\
\rf{7}  & O sistema deve extrair texto de documentos digitalizados por
\textsc{ocr}. & Should \\
\rf{8}  & O sistema deve construir uma linha temporal cronológica,
resolvendo datas incompletas. & Should \\
\rf{9}  & O sistema deve gerar narrativas a partir dos factos, permitindo
escolher tema, tom e idioma (PT/EN). & Must \\
\rf{10} & O sistema deve gerar um vídeo documental narrado a partir de uma
narrativa. & Must \\
\rf{11} & O sistema deve permitir organizar fotografias, histórias e
vídeos em projetos (coleções curadas). & Should \\
\rf{12} & O sistema deve executar gerações longas em segundo plano, com
acompanhamento de estado e notificações. & Should \\
\rf{13} & O sistema deve permitir descarregar e reproduzir os vídeos
gerados. & Must \\
\rf{14} & A interface deve estar disponível em português europeu
(omissão) e inglês. & Could \\
\rf{15} & O sistema deve expor o estado de saúde das suas dependências
(sondas de \emph{health}). & Could \\
\rf{16} & O sistema deve permitir gerir as histórias geradas --- editar,
regenerar com \emph{feedback}, marcar como favoritas, pesquisar/filtrar e
exportar para PDF. & Should \\
\rf{17} & O sistema deve permitir visualizar e editar a árvore familiar de
forma interativa (pesquisar pessoas, destacar linhagem, adicionar
familiares e exportar a árvore como imagem). & Should \\
\rf{18} & O sistema deve permitir escolher a voz da narração e incluir uma
faixa de legendas sincronizada e comutável no vídeo. & Could \\
\rf{19} & O sistema deve permitir etiquetar as pessoas presentes em cada
fotografia e gerir a galeria de cada pessoa. & Should \\
\bottomrule
\end{longtable}

\section{Requisitos Não Funcionais}
\label{sec:rnf}

\begin{longtable}{@{}p{1.6cm}p{3.4cm}p{8.6cm}@{}}
\caption{Requisitos não funcionais do sistema.}\label{tab:rnf}\\
\toprule
\textbf{ID} & \textbf{Atributo} & \textbf{Descrição} \\
\midrule
\endfirsthead
\multicolumn{3}{c}{\tablename\ \thetable\ -- continuação}\\
\toprule
\textbf{ID} & \textbf{Atributo} & \textbf{Descrição} \\
\midrule
\endhead
\rnf{1} & Segurança & Validação criptográfica de tokens; isolamento estrito
entre inquilinos; proteção contra abuso por \emph{rate limiting};
validação de \emph{uploads}. \\
\rnf{2} & Privacidade & Possibilidade de inferência local, sem envio de
dados do arquivo para serviços de terceiros; conformidade com os
princípios do RGPD. \\
\rnf{3} & Robustez & Degradação graciosa na ausência de dependências
opcionais (base vetorial, \textsc{llm} local, fila de tarefas). \\
\rnf{4} & Observabilidade & \emph{Logging} estruturado e sondas de saúde
profundas sobre todas as dependências. \\
\rnf{5} & Desempenho percebido & \emph{Cache} de estado de servidor,
\emph{polling} controlado e \emph{feedback} imediato ao utilizador. \\
\rnf{6} & Manutenibilidade & Tipagem estática (Python e TypeScript),
modularidade e cobertura de testes automatizados. \\
\rnf{7} & Portabilidade & Empacotamento em contentor Docker com todas as
dependências de sistema. \\
\rnf{8} & Usabilidade & Interface intuitiva, responsiva (\emph{desktop} e
móvel), com tema claro/escuro e suporte a acessibilidade. \\
\rnf{9} & Custo & Operação assente em ferramentas de código aberto e
camadas gratuitas (custo $\approx 0$). \\
\bottomrule
\end{longtable}

\section{Atores e Casos de Uso}
\label{sec:casos_uso}

O sistema tem um ator principal — o \textbf{Utilizador autenticado}
(o ``dono'' de um arquivo familiar) — e dois atores secundários, que são
sistemas externos: o \textbf{Serviço de IA} (Groq/Ollama/Gemini) e o
\textbf{Supabase} (autenticação, base de dados e armazenamento). A
\Cref{fig:casos_uso} apresenta o diagrama de casos de uso.

\begin{figure}[H]
\centering
% \includegraphics[width=0.92\linewidth]{casos_uso.png}
\figph[7.5cm]{Diagrama UML de casos de uso. Ator ``Utilizador''
ligado aos casos: Autenticar-se; Carregar fotografias; Importar GEDCOM;
Gerar narrativa; Gerar vídeo; Gerir projetos; Consultar tarefas; Gerir
definições. Atores secundários ``Serviço de IA'' e ``Supabase'' ligados
aos casos de geração e de persistência. Sugestão: desenhar em
draw.io / PlantUML.}
\caption{Diagrama de casos de uso do sistema.}
\label{fig:casos_uso}
\end{figure}

\section{\emph{User Stories} Principais}
\label{sec:user_stories}

As \emph{user stories} seguintes capturam a intenção do utilizador e
ligam-se aos requisitos funcionais correspondentes:

\begin{itemize}[leftmargin=*]
  \item \textbf{US1} — \emph{``Como utilizador, quero carregar as
  fotografias antigas da minha família arrastando-as para o ecrã, para
  não ter de as enviar uma a uma.''} (\rf{3})
  \item \textbf{US2} — \emph{``Como utilizador, quero importar a árvore
  genealógica que já tenho em GEDCOM, para que as narrativas conheçam as
  relações familiares.''} (\rf{4})
  \item \textbf{US3} — \emph{``Como utilizador, quero pedir uma história
  sobre um tema específico (por exemplo, a emigração do meu avô),
  escolhendo o tom, para obter uma narrativa fiel ao que tenho em
  mente.''} (\rf{9})
  \item \textbf{US4} — \emph{``Como utilizador, quero transformar uma
  história num vídeo documental narrado, para a poder partilhar e
  preservar.''} (\rf{10})
  \item \textbf{US5} — \emph{``Como utilizador, quero organizar o meu
  material em projetos temáticos, para separar diferentes ramos ou
  ocasiões da família.''} (\rf{11})
  \item \textbf{US6} — \emph{``Como utilizador, quero que as gerações
  longas corram em segundo plano e me avisem quando terminam, para não
  ficar à espera.''} (\rf{12})
  \item \textbf{US7} — \emph{``Como utilizador preocupado com a
  privacidade, quero que os meus dados não saiam da máquina sempre que
  possível.''} (\rnf{2})
\end{itemize}
